\section{Fazit}

Die Umsetzung verdeutlicht sowohl die Vorteile als auch die Nachteile einer Microservice-Architektur.
Die fachliche Aufteilung, klar definierte Schnittstellen und
getrennte Datenbanken ermöglichen eine weitgehend unabhängige Entwicklung der einzelnen Services.
Durch die Kombination aus synchroner und asynchroner Kommunikation konnten unterschiedliche Formen der Servicekopplung praktisch umgesetzt und erprobt werden.

Gleichzeitig erhöht die Verteilung der Anwendung den technischen Aufwand.
Betrieb und Fehlersuche erstrecken sich über mehrere Services, Container und Kommunikationswege.
Serviceübergreifende Vorgänge können nicht durch eine gemeinsame Datenbanktransaktion abgesichert werden und
benötigen deshalb zusätzliche Kompensationslogik.
Zudem bleibt der Kaufvorgang von synchron aufgerufenen Services wie dem Product Service und dem User Service abhängig.

Für größere Systeme mit mehreren Entwicklungsteams,
unterschiedlich belasteten Teilbereichen oder
hohen Anforderungen an die Fehlerisolation kann eine Microservice-Architektur sinnvoll sein.
Für den begrenzten Umfang der umgesetzten Beispielanwendung wäre ein modularer Monolith einfacher zu entwickeln und
zu betreiben gewesen.
Als Lernprojekt war die gewählte Architektur dennoch geeignet, da ihre Auswirkungen auf Implementierung,
Kommunikation, Datenhaltung und Fehlerbehandlung praktisch untersucht werden konnten.